\begin{abstract}
视觉状态空间模型必须先把图像排成序列，再执行选择性扫描，因此路线决定哪些信息先进入历史、哪些信息靠近最终读出。多数视觉路线对所有图像固定。本文提出\method：用确定性显著性场为每张图像生成路线，再沿路线直接读取原图像素。完整等高线树提供延续、分裂、汇合、源点和终止关系；只有值得扫描的常规轮廓才插入稀疏执行树。每个保留轮廓先按周长确定点数，再由均匀试探点读取同一张Sobel场：法线方向响应的平均符号为整圈选择低显著性普通读取侧，平滑后的响应强度重分配固定点数，最大强度位置确定共同起止点 $p_0$。Normal Mamba从首次碰撞点读回带Tag的轮廓点。两次共享参数的Contour Mamba扫描从 $p_0$ 出发，按相反方向遍历同一组闭环点，并在带Tag的 $p_0$ 副本结束。每条树枝的终止圈读取两侧法线，再用第二对闭环汇集。所有轮廓独立编码完成后，Tree Mamba才从低显著性向高显著性汇总带Tag输出：分裂复制摘要，汇合采用与输入顺序无关、考虑采样支撑量的融合。最后，各终止叶按实际读取区域的显著性排序，并由最终Mamba完成分类。初版路线完全确定，不使用CNN特征Stem。本文完整规定路线、Token接口、边界情况、默认值和结构性质，不作实证性能声明。
\end{abstract}
